% ---------------------------------------------------------------------------
%  04-parallel.tex -- the OpenMP architecture proper
% ---------------------------------------------------------------------------
\section{The parallel architecture}
\label{sec:parallel}

\FigLayers

The engine is one function, \code{reach\_engine} in \file{src/netreach.cpp},
parameterised by the target predicate and by whether subsumption must be exact.
Every query -- \code{E<>}, \code{A[]}, \code{E<> deadlock} -- runs on it, so
there is one parallel implementation to get right rather than one per query
form.

\subsection{Level-synchronous BFS}

The search is breadth-first with an explicit barrier at each level
(\cref{fig:level}).  A level does four things:

\begin{enumerate}[leftmargin=1.7em, itemsep=2pt]
\item \textsc{expand} -- in parallel, expand every frontier state and insert
      its successors directly into the shared passed set;
\item \textsc{collect} -- in parallel, gather the states that survived
      insertion at this level into the next frontier;
\item \textsc{compact} -- in parallel, reclaim tombstones inside each stripe;
\item \textsc{gc} -- occasionally, mark-and-sweep the zone interner.
\end{enumerate}

\FigLevel

Level synchronisation costs some idle time at each barrier -- the ``tail idle''
in \cref{fig:level} -- and a work-stealing DFS would avoid it.  It is kept
deliberately, for a reason specific to \emph{subsumption-based} search: BFS
order minimises the number of stored states, because a zone reached by a
shorter path tends to be more general and therefore kills more of its
successors' siblings.  Trading a few percent of idle time for a smaller state
space is the right trade when the state space is the bottleneck, which it is on
every model here.

\begin{intuition}
The workers go through the search in waves.  Everyone expands the current wave
of states at the same time; when everyone is done, the survivors become the
next wave.  Waiting at the end of each wave wastes a little time, but exploring
in strict wave order means you find each state by its shortest path, and short
paths give more general shapes, and more general shapes let you throw away more
work later.
\end{intuition}

\subsection{Fused expansion and insertion}

\begin{center}
\begin{minipage}{0.93\linewidth}
\begin{lstlisting}
#pragma omp parallel num_threads(nthreads)
{
  const int tid = omp_get_thread_num();
  std::vector<SymState> buf;      // reused across frontier states
  SymState cur;                   // reused: reconstructed frontier state
#pragma omp for schedule(dynamic, expand_chunk) nowait
  for (long i = 0; i < (long)frontier.size(); ++i) {
    if (found.load(std::memory_order_relaxed)) continue;
    decode_key(frontier[i].first, cur.locs.data(), cur.ints.data());
    const size_t ns = gen_succ_reuse(net, cur, zi.get(frontier[i].second), ...);
    for (size_t k = 0; k < ns; ++k) { /* target test, then insert */ }
  }
}
\end{lstlisting}
\end{minipage}
\end{center}

Three things in that loop are deliberate.

\paragraph{No routing barrier.}
A worker inserts each successor into the shared table as soon as it produces
it.  The alternative -- generate into per-thread buffers, then route them to
owners at a barrier -- introduces a serial phase per level whose cost grows
with level width, which is precisely what Amdahl's law punishes.  Fusing
expansion and insertion removes that phase entirely.

\paragraph{Dynamic scheduling with small chunks.}
Frontier states differ wildly in cost: one may fire two edges, another may fire
a broadcast synchronisation with a combinatorial number of receiver
combinations.  A static schedule leaves lanes idle; \code{schedule(dynamic, 8)}
keeps them within a chunk of each other.  The chunk size is a CLI knob
(\code{--omp-chunk}).

\paragraph{Successor storage is reused, not allocated.}
\code{gen\_succ\_reuse} writes each successor into a slot of a per-thread buffer
that survives across frontier states, and returns a count; the caller reads
\code{buf[0..n)}.  After a short warm-up, generating a successor performs
\emph{zero} allocations.  This matters far more in parallel than sequentially:
\code{malloc} holds a process-global lock, so per-successor allocation makes
all workers serialise on the allocator.  What allocation remains goes through a
per-thread free-list allocator (\file{dbm\_pool.hpp}) rather than the system
heap.

\subsection{The striped passed set}

\FigStripes

The passed set is $\NB \ge \max(16\,384,\, 64 \cdot \text{threads})$ stripes,
each an independent flat open-addressed table with its own reader--writer lock.
A discrete key selects its stripe by the high bits of a multiplicative hash.

This is \ltsmin{}'s shape -- one shared table with fine-grained locking and no
static partition of the state space -- rather than the older
``partition-and-route'' shape.  The difference is not cosmetic: with a static
partition, a model whose states cluster in a few partitions starves most
workers, and every level needs a routing step.  With a shared table, load
balance follows the hash, and there is nothing to route.

The locks themselves are a custom spin-then-yield reader--writer lock
(\code{SpinRW}, $8$ bytes).  With \code{std::shared\_mutex} each collision cost
$0.3$--$1.8$ \emph{milliseconds}, because the loser blocks and waits for a
scheduling quantum to service a critical section that runs in microseconds; on
\code{phil13} that convoy consumed $378$ of $759$ CPU-seconds.  Spinning turns a
collision into a few hundred nanoseconds; the bounded spin falls back to
\code{yield()} so that oversubscription (threads $>$ cores) degrades gracefully
instead of burning quanta against a descheduled holder.

\subsection{The two-phase lock protocol}

\FigProtocol

This is the centre of the design.  Insertion is split into a read-only prefix
and a mutating remainder:

\begin{center}
\begin{minipage}{0.93\linewidth}
\begin{lstlisting}
smtx[si].lock_shared();
const bool covered = smap[si].probe_kill(dk, s, exact_sub, state_bounds);
smtx[si].unlock_shared();
if (covered) continue;                 // 83% of attempts end here

smtx[si].lock();                       // only real inserts get here
if (smap[si].add(std::move(s), dk, level, exact_sub, state_bounds, tid) >= 0)
  touched.insert(dk);
smtx[si].unlock();
\end{lstlisting}
\end{minipage}
\end{center}

\code{probe\_kill} is \code{add}'s fast path verbatim, with everything that
writes removed.  It answers ``would \code{add} have returned \emph{drop}?''.
Because shared locks do not exclude one another, the $83\%$ of attempts it
disposes of never serialise.

Two correctness obligations arise, and they are asymmetric:

\begin{itemize}[leftmargin=1.2em, itemsep=2pt]
\item A \textbf{false ``no''} (probe says not covered, but by the time we take
      the exclusive lock it is) is harmless: \code{add} redoes the whole test
      and will drop the state.  It costs one wasted lock acquisition.

\item A \textbf{false ``yes''} would be a wrong answer: we would drop a state
      that is not in fact covered, and the search could miss reachable
      behaviour.
\end{itemize}

So the entire protocol rests on: \emph{a ``yes'' cannot go stale}.  That is not
obvious --- between the two phases, another worker may retire the very zone
that answered ``yes''.  \cref{thm:coverage} and \cref{cor:probe} prove it does
not matter, because retirement in this data structure always replaces a zone
by something containing it.

\begin{intuition}
Two workers, one shared shelf of shapes.  Worker~A looks at the shelf without
locking anybody out and concludes ``my shape is already covered, I can skip
it''.  Meanwhile worker~B removes the very shape that covered it.  Is A's
conclusion now wrong?  No --- because B is only allowed to remove a shape by
putting a \emph{bigger} one in its place.  The shelf's coverage only ever
grows.  That single structural property is what makes the lock-free read legal.
\end{intuition}

\subsection{Bucketed frontier collection}

After the barrier, each worker holds a set of discrete keys it touched, and
those sets must be unioned so that each key is scanned exactly once when
building the next frontier.  Doing that with one shared \code{unordered\_set}
was the largest remaining serial phase ($0.149$\,s of a $3.83$\,s
\code{arbiter16} run at $8$ threads) and it grows with level width.

Instead, keys are bucketed by hash into $\ge 4 \cdot \text{threads}$ buckets and
each bucket is deduplicated independently, in parallel, with the
\code{for\_level} scan fused into the same loop.  Since a key hashes to exactly
one bucket, that bucket sees every copy of it, and the deduplicated key set is
identical to the serial union's (\cref{lem:bucket}).  The buckets are taken
from the same hash as the stripe index, so a bucket's scans stay inside a
contiguous slice of stripes instead of scattering over all $\NB$ tables.

\subsection{Early termination}

A worker that finds a target state sets a relaxed atomic flag; other workers
observe it at the top of their next iteration and skip the remaining work.  The
flag is read \emph{relaxed} on purpose --- it is a hint, not a synchronisation
point; the authoritative read happens after the barrier.  Delaying the exit by
a few iterations costs nothing, whereas an acquire load on every frontier state
would cost on every iteration of the hottest loop.

(The single-automaton path in \file{src/reach.cpp} uses OpenMP's
\code{cancel for} / \code{cancellation point for} instead, which stops workers
at their next iteration boundary rather than at the end of their chunk.  Both
are correct; the network engine's frontier iterations are coarse enough that
the flag is sufficient.)

\subsection{Bounded memory: compaction and garbage collection}

Subsumption and merging mark nodes dead but do not free them, and interning
accumulates records for zones that no longer have a referrer.  Two reclamation
passes run between levels, both parallel over stripes, and both with
\emph{amortised} triggers:

\begin{itemize}[leftmargin=1.2em, itemsep=2pt]
\item \textbf{Stripe compaction} rebuilds a stripe's chains live-only once a
      third of the stripe is tombstones.  A rebuild costs $O(\text{live})$ and
      drops \emph{every} tombstone, so the work is amortised $O(1)$ per
      tombstone.  Trigger and floor are both per stripe: an earlier global-scale
      floor meant compaction never fired on any stripe, and $54\%$ of the arena
      on \code{phil13} was garbage that every chain walk had to step over.

\item \textbf{Interner GC} marks every zone index still referenced by a live
      node or by the next frontier, then sweeps the rest onto free lists.  The
      trigger is \emph{growth-based}: sweep once the allocated-zone count has
      doubled since the last sweep.  The earlier per-level trigger ran a full
      $O(\text{arena})$ mark-and-sweep on every level and was $30\%$ of a
      \code{arbiter17} run --- serial work growing with the state space while
      the parallel part did not.  \cref{prop:gc} states the amortisation and
      the memory bound it trades for.
\end{itemize}

The mark phase itself is parallel: every writer stores the same value, and
nothing reads the array until the sweep, so concurrent marks cannot disagree.
The stores are nevertheless relaxed atomics, because concurrent plain stores to
one byte are formally a data race; on x86 and ARM a relaxed byte store is the
same instruction as a plain one, so the parallelism is free.
